% Abstract: one paragraph, no subheadings, ordered motivation and aim, methods, results, conclusion; at most 250
% words (counted on the rendered text).
\begin{abstract}
\noindent General-purpose large language models (LLMs) can write and run code, but whether they can use it to classify
raw electrocardiograms (ECGs) without ECG-specific training has, to our knowledge, not been evaluated. We gave one
such LLM (\ModelID) a Python sandbox and, for each raw 12-lead PTB-XL recording, a file with the signal,
the patient's age and sex, and a prompt. The prompt asked the model to filter the signal, form median beats and print
ECG measurements, then apply listed diagnostic criteria and give a probability for each of the five diagnostic
superclasses; it contained no labelled examples. The prompt was developed on the validation fold and frozen. On the
benchmark test fold ($n{=}\NTest$), the model wrote and ran code \CodeRunsMin--\CodeRunsMax{} times per record
(\PHSandboxCodeRunsTotal{} code runs). The macro area under the receiver operating characteristic curve was
\MainAUROC{} (95\% CI \MainAUROCLow--\MainAUROCHigh), compared with \AgeSexAUROC{} for a logistic regression using
only age and sex. On the same benchmark, the best published supervised and zero-shot ECG--language models reach
\BestPublishedAUROC{} and \BestZeroShotAUROC, respectively. A post hoc analysis of the transcripts suggested that most
ECG intervals and axes the model stated were read by the model from median-beat samples its code printed, not
computed. Run again offline, the model's code reproduced the sandbox's output, up to rounding and other trivial
differences. A general-purpose LLM writing its own code thus scored between published zero-shot and supervised models,
and every answer can be traced to the code, its output and the measurements the model stated.
\end{abstract}

\noindent\textbf{Keywords:} electrocardiography; large language models; code generation; PTB-XL; zero-shot
classification; interpretability
